%% kapa: exam=1 round=33 year=2022 period=2 form=A subject=accounting subject_ko=회계학 seq=32 number=72
%% source: Q-Net 감정평가사 기출문제 (cst00309) artlSeq=5209997 file=raw/1008/5209997/2022년도 제33회 감정평가사 1차시험 2교시(감정평가 관계법규, 회계학).pdf
%% conversion: kordoc 4.24.51 -> tools/gen1.py (auto); review=proofread-sonnet; flags=none
\begin{question}{72}
원가에 관한 설명으로 옳지 않은 것은?
\begin{choices}
  \item 가공원가(전환원가)는 직접노무원가와 제조간접원가를 합한 금액이다.
  \item 연간 발생할 것으로 기대되는 총변동원가는 관련범위 내에서 일정하다.
  \item 당기제품제조원가는 당기에 완성되어 제품으로 대체된 완성품의 제조원가이다.
  \item 기초고정원가는 현재의 조업도 수준을 유지하는데 기본적으로 발생하는 고정원가이다.
  \item 회피가능원가는 특정한 의사결정에 의하여 원가의 발생을 회피할 수 있는 원가로서 의사결정과 관련있는 원가이다.
\end{choices}
\end{question}
\answer{2}
